% TIME 2027 — Extended Abstract submission
%
% Target venue: 34th International Symposium on Temporal Representation
% and Reasoning (TIME 2027); venue and dates TBC.
% Category: Extended abstract (4 pages, excluding references and appendix).
% Format:   OASIcs (oasics-v2021).
%
% Build:
%   1. Download oasics-v2021 from
%      https://submission.dagstuhl.de/documentation/authors
%      and place oasics-v2021.cls (plus its supporting files) in this
%      directory.
%   2. pdflatex extended-abstract && pdflatex extended-abstract
%
% This file uses an inline thebibliography environment so no .bib is
% required at draft stage. Switch to BibTeX (plainurl) before final
% submission.

\documentclass[a4paper,USenglish,cleveref,autoref,thm-restate]{oasics-v2021}

\title{Time as Intervals, Not Instants: \texorpdfstring{\\}{ }%
  A Practical Implementation of Allen's Interval Algebra Over ISO~8601}
\titlerunning{Time as Intervals, Not Instants}

\author{Kip Cole}{Independent}{kipcole9@gmail.com}{}{}

\authorrunning{K.~Cole}
\Copyright{Kip Cole}

\ccsdesc[500]{Information systems~Temporal data}
\ccsdesc[300]{Software and its engineering~Domain specific languages}
\ccsdesc[300]{Theory of computation~Modal and temporal logics}

\keywords{Temporal intervals, Allen's interval algebra, ISO 8601,
  calendrical systems, partial dates, set operations on time,
  domain-specific languages}

\relatedversion{Source repository: \url{https://github.com/elixir-tempo/tempo}}

% The inline thebibliography below uses \doi{...} the way bibtex's
% `plainurl` style would emit it. The OASIcs class defines \doi for
% the article's own header but not for bibitem entries, so provide
% it here. Switch to BibTeX (plainurl) before final submission and
% this can be removed.
\providecommand{\doi}[1]{\href{https://doi.org/#1}{doi:\texttt{#1}}}

\begin{document}

\maketitle

\begin{abstract}
  Mainstream programming languages model temporal values as scalar
  instants: a \texttt{Date} is a point, a \texttt{DateTime} is a point,
  and a partial value such as ``June 2026'' has no canonical
  representation at all. Yet the temporal values that humans, calendars,
  and business logic actually manipulate are spans on the timeline ---
  the year 2026, the afternoon of the 15th, an archaeological period,
  a meeting. Tempo is an Elixir library that reverses the polarity.
  Every expression --- a value, an interval, a recurrence, a set ---
  denotes an ordered sequence of bounded half-open spans; a value
  written to a resolution denotes one unit of it; set-theoretic
  combination is closed over the one type; and comparison is defined by
  Allen's thirteen relations rather than by scalar ordering. A written
  interval is read as it is said aloud: an end written to a calendar
  unit is included, an end written to a clock unit is a boundary. The
  result eliminates a class of off-by-one and ``end-of-day'' bugs by
  construction, gives partial ISO~8601 values (\texttt{2026},
  \texttt{2026-06}, \texttt{2026-06-XX}) an unambiguous denotational
  semantics, and lifts Allen's algebra from a textbook abstraction to an
  everyday predicate vocabulary (\texttt{overlaps?}, \texttt{within?},
  \texttt{adjacent?}). This extended abstract describes the model as
  nine laws, the frames on which spans are placed, the reading of
  written ends, and the design of set operations over spans and sets
  of them, and reports on the implementation written to the laws.
\end{abstract}

\section{Motivation: the instant--span impedance mismatch}

Every mainstream temporal library --- Python's \texttt{datetime}, Java's
\texttt{java.time}, JavaScript's \texttt{Temporal}, Elixir's
own \texttt{Date}/\texttt{Time}/\texttt{DateTime} --- shares a design
choice that goes largely unexamined: a temporal value is a
\emph{scalar instant} on the timeline. A day is the midnight at its
start, a month is the midnight on its first, a year is the midnight
on January~1. The span that humans actually mean when they say ``June
2026'' is materialised, if at all, by ad-hoc helpers
(\texttt{start\_of\_month}, \texttt{end\_of\_day}) that each project
re-implements with its own conventions about inclusivity, time zones,
and DST.

This instant-first model has three persistent consequences:

\paragraph*{Ambiguity at the type level.} The expression
\texttt{Date.new(2026, 6, 15)} denotes a 24-hour span in business
prose and an instant in the type system. The disagreement is resolved
informally, per call site, and is a recurring source of off-by-one
bugs at month, year, and DST boundaries. Tiwari et al.'s recent
empirical study of 151 date/time bugs in open-source Python projects
identifies time-zone-related mistakes as the largest single category,
with most bugs traceable to misconceptions about library API
conventions and edge-case behaviour~\cite{tiwari2025time}; Liva's
earlier work on Java programs reaches a structurally similar
conclusion~\cite{liva2018}.

\paragraph*{No representation for partial dates.} A user who writes
``2022'' or ``the 1560s'' or ``approximately June 2026'' has produced
a perfectly meaningful temporal value, but the type system has no
home for it. Solutions range from sentinel values
(\texttt{Date.new(2022, 1, 1)} stands in for ``the year'') to parallel
type hierarchies (EDTF~\cite{edtf}) that exist outside the standard
library, or to dedicated calendric type systems such as Bry and
Spranger's CaTTS for Semantic Web query
languages~\cite{bry-spranger-multicalendar}. The fragmentation is
structural, not incidental.

\paragraph*{Comparison is reduced to scalar ordering.} Two dates are
either \texttt{<}, \texttt{=}, or \texttt{>}. The richer relational
structure that Allen identified in 1983~\cite{allen1983} ---
\textit{precedes}, \textit{meets}, \textit{overlaps}, \textit{starts},
\textit{during}, \textit{finishes}, \textit{equals}, and their inverses
--- is invisible at the type level, and reimplemented case-by-case in
business logic that interrogates whether a meeting ``overlaps'' the
workday or a contract ``contains'' a billing period.

Tempo's design is the observation that all three problems dissolve
simultaneously if intervals are made the primitive type. There are no
points: an apparent ``instant'' is the interval of one unit at the
finest declared resolution, in the same spirit as the time ontology
$T_{bounded\_meeting}$ of Gr\"uninger and
Li~\cite{gruninger2017allenontology}, in which intervals are the only
entities in the domain and Allen's relations are derivable from a
single primitive (\textit{meets}).

\paragraph*{Empirical coverage.} Against the 151 Python bugs
catalogued by Tiwari et al., the design choices described in the
remainder of this paper address every in-scope bug class: 43 (28\%)
are \emph{avoided by construction} (the na\"ive / aware datetime
split has no Tempo analogue), 56 (37\%) are \emph{directly addressed}
by explicit machinery (zone propagation through set operations,
refusal of a reading the clock skips and IXDTF-offset fold
disambiguation per RFC~9557~\S 4.5, calendar-aware arithmetic,
half-open materialisation, microsecond-precision durations), 28 (19\%)
are partial, and 24 (16\%) are out of scope (Python deprecations and
third-party-library logic such as holiday tables). \emph{Zero} bugs
are classified as architectural gaps. The classification was made
against the library's first implementation; the rewrite described in
Section~\ref{sec:implementation} kept every mechanism it relies on.
The per-bug classification is provided as supplementary material at
the repository.

\section{Expressions, spans and frames}
\label{sec:model}

Tempo's model is nine laws, and the implementation is written to them
and to nothing else. The first is the one from which the rest follow:
an \emph{expression} --- a value such as \texttt{2026-06}, an interval
such as \texttt{2026-06-01/2026-06-20}, a recurrence such as
\texttt{R5/2026-06-15/P1D}, a set of them --- \emph{denotes an ordered
sequence of half-open spans}, $\mathit{spans}(E)$, possibly empty,
possibly infinite, and every operation on an expression is defined on
$\mathit{spans}(E)$. A duration denotes a length and is not an
expression. One data structure holds every expression, with its form
as a tagged sum (a value, a span, a recurrence, an all-of set, a one-of
set, a conditional member), its calendar, its zone and the caller's
metadata; there is no separate interval type and no separate set type.

A span is $[\mathit{first}, \mathit{last})$, a pair of positions.
Internally this is the standard finite model of $T_{bounded\_meeting}$
characterised by Gr\"uninger and Li~\cite{gruninger2017allenontology}
(their Theorem~15): vertices of the directed graph
$\mathfrak{M}^{bounded\_meeting}$ correspond to intervals, edges to the
\textit{meets} relation, and structures up to isomorphism faithfully
capture the ontology's models.

\paragraph*{Resolution is meaning.} A value written to a resolution
denotes one unit of it, $[v, \mathit{next}_r(v))$:
\begin{center}
\begin{tabular}{lcl}
  \texttt{2026}              &$\equiv$& $[\texttt{2026-01-01},\, \texttt{2027-01-01})$ \\
  \texttt{2026-01}           &$\equiv$& $[\texttt{2026-01-01},\, \texttt{2026-02-01})$ \\
  \texttt{2026-01-15}        &$\equiv$& $[\texttt{2026-01-15},\, \texttt{2026-01-16})$ \\
  \texttt{2026-01-15T10}     &$\equiv$& $[\texttt{2026-01-15T10:00},\, \texttt{2026-01-15T11:00})$ \\
\end{tabular}
\end{center}
The partial-date problem disappears: \texttt{2026} is not an
under-specified instant, it is a definite interval of one year (the
distinction is semantic, not a claim about measurement precision or
accuracy). The ``last day of the month'' problem disappears: the
upper bound of \texttt{2026-01} is
\texttt{2026-02-01}, computed from the calendar arithmetic of the
target calendar (Gregorian, Hebrew, Islamic Civil, \dots), and is
correct by construction in leap years and across month-length
variation. The written units are kept as written; nothing drills into
them, and a value and the interval of the same extent are different
expressions that enumerate differently (a month enumerates its days;
the interval \texttt{2026-01/2026-01} enumerates months, one of
them).

A note on terminology: \emph{resolution} as used
throughout this paper is a property of an input value
(\texttt{2026} carries year-resolution, \texttt{2026-06-15} carries
day-resolution) used for materialisation, default
iteration, and display. It is \emph{not} the granularity of
Bettini, Wang and Jajodia~\cite{bettini-granularities} (a lattice
of partitions of the time line that participates in Allen
comparisons), and it does \emph{not} appear in
$T_{bounded\_meeting}$, where intervals are individuated solely by
their position in the \textit{meets}-structure. Once an interval is
materialised to its $[\mathit{first}, \mathit{last})$ form,
resolution plays no further role in Allen comparisons or in
set-theoretic operations. The word is also distinct from
\emph{precision} and \emph{accuracy} in the measurement-science
sense: Tempo makes no claim about the precision of a value's
construction, only about the extent it denotes.

\paragraph*{Written ends are read as they are said.} An interval
written with two ends, \texttt{a/b} in ISO~8601, runs from the start
of \texttt{a}. Where it runs to depends on what \texttt{b} names. A
date names a period of the calendar, and ``2026 to 2028'' and ``1 to
3 June'' include their last day in speech, in EDTF (\texttt{1964/2008}
is ``beginning sometime in 1964 and ending sometime in 2008'') and in
ISO~8601-2 Annex~A. A time of day names a reading of the clock, and a
reading is a boundary: ``nine to five'' ends when the clock shows five,
not a minute later. So an end written to a \emph{calendar} unit (a
year, a month, a week, a date) is read \emph{through}, and an end
written to a \emph{clock} unit (an hour, a minute, a second) is read
as a boundary:
\begin{center}
\begin{tabular}{lll}
  \texttt{2026/2028}                     & $[\texttt{2026-01-01},\, \texttt{2029-01-01})$ & three years \\
  \texttt{2026-06-01/2026-06-03}         & $[\texttt{06-01},\, \texttt{06-04})$            & three days \\
  \texttt{T09/T17}                       & $[\texttt{09:00},\, \texttt{17:00})$            & eight hours \\
  \texttt{2026-06-15/2026-06-15T17}      & $[\texttt{06-15T00:00},\, \texttt{06-15T17:00})$ & seventeen hours \\
  \texttt{2026-06-15T10/2026-06-16}      & from 10:00 on the 15th through the 16th & thirty-eight hours \\
\end{tabular}
\end{center}
The reader makes the half-open bound and the writer writes a span back
as it was read, so text round-trips. The rule needs nothing of a
calendar, only the class of the unit, and it lets the half-open model
stand without a notation for an exclusive end, which ISO~8601 does not
have. It is the one semantic choice in the library with no direct
precedent in the formalisms it draws on, and Section~\ref{sec:discussion}
puts it to the symposium.

\paragraph*{Frames.} A span lies on one of three frames. A value with a
year and a zone, or a written offset, lies on the \emph{UTC line}; a
value with a year and neither lies on the \emph{floating line}; a value
with no year (\texttt{T10:30}, \texttt{--06-15}, a Monday) lies on a
\emph{cycle}, the day for a time of day, the week for a weekday, the
year for a month and day. A position is an integer: a microsecond tick
from a fixed origin on a line, a tick from the cycle's start on a
cycle. Two spans relate, compare or combine only on one frame; a
floating value beside a zoned one, or a time of day beside a date, is a
named error and never a guess. A value on a cycle is, in the ontology's
terms, not an individual but an open expression with the year free; the
\texttt{:within} option, and the operations \texttt{at/2},
\texttt{on/2} and \texttt{in\_zone/2}, bind it, placing it on each
period of a window, and the result is the intervalised individual
$T_{bounded\_meeting}$ talks about. Set operations on a cycle cut a
span that runs through the cycle's end (\texttt{T22/T02}) for the
sweep and join it after, so the algebra is the same on all three
frames. The ontology is silent on schemas; frames, like resolution,
sit in the engineering layer above it.

\paragraph*{Half-open by design.} The half-open $[\mathit{first},
\mathit{last})$ convention is not notational shorthand; it is enforced
at every level of the API. Adjacent intervals concatenate cleanly:
$[a, b) \cup [b, c) = [a, c)$ with neither overlap nor gap. Iteration
yields a sequence whose union is exactly the parent interval. No span
has zero extent, in line with $T_{bounded\_meeting}$'s exclusion of
degenerate intervals from the domain. Library code that treats the
stored upper bound as inclusive is, by the design contract, a defect.
Alternative timestamp representations have been explored at TIME ---
notably Dyreson and Ahsan's log-segmented
timestamps~\cite{dyreson2021logsegmented,dyreson2023logsegmented},
which decompose a period into power-of-two segments so that standard
B$^+$-tree indexes can support sequenced and nonsequenced queries on
an unmodified DBMS. The two designs occupy different points in the
same space: log-segmented timestamps optimise query evaluation in a
relational engine; Tempo's half-open spans optimise expressiveness
and uniformity at the application layer.

\section{Allen's algebra as the predicate vocabulary}

Tempo exposes Allen's thirteen relations~\cite{allen1983} directly as
the result of its core comparison primitive. The first-order ontology
underlying the algebra is characterised by Gr\"uninger and
Li~\cite{gruninger2017allenontology}, who identify the time ontology
logically synonymous with Allen's interval algebra and prove a
representation theorem for its models; Tempo is, in effect, an
ergonomic software realisation of that ontology over the
calendar-aware datetime values defined by ISO~8601. The thirteen
relations are returned directly from the relation primitive:
\begin{verbatim}
  Tempo.relation(a, b) ::
    :precedes | :meets | :overlaps | :finished_by |
    :contains | :starts | :equals |
    :started_by | :during | :finishes | :overlapped_by |
    :met_by | :preceded_by
\end{verbatim}
The relation is read from four integer comparisons of the two spans'
positions (start against start, end against end, and each start
against the other's end); every pair of spans stands in exactly one
relation.

Beyond the relations themselves, the library ships Allen's
\emph{composition} table --- the transitivity result of~\cite{allen1983}
Fig.~4 --- as a first-class operation, \texttt{Tempo.Allen.compose/2}:
given $A\;r_1\;B$ and $B\;r_2\;C$, it returns the set of relations that
can hold between $A$ and $C$. The 169 cells are not transcribed: they
are derived when the library compiles, by placing three intervals every
way on eight points (six endpoints need at most six distinct values, so
the enumeration is complete) and classifying each pair with the same
four comparisons, and then held to the published table by test. The
classifier and the composition table therefore cannot disagree.

Composition drives a propagation algorithm. Lifting it from single
relations to \emph{sets} of them gives the operations a qualitative
constraint network is built from --- the converse of a set, the
narrowing of one set by another, and set composition --- and Allen's
path-consistency step is then their composite,
$R(i,k) \leftarrow R(i,k) \cap (R(i,j) \circ R(j,k))$, iterated to a
fixpoint. This is what lets a companion package, \texttt{tempo\_network},
reason about intervals that are \emph{not} grounded: a period known
only as ``before or after the rebuild, never during'' is the set
$\{\textit{precedes}, \textit{preceded\_by}\}$, and propagation derives
what that implies elsewhere in the network. Pruning is sound --- every
relation removed is genuinely impossible --- but path consistency does
not decide satisfiability, which is NP-complete for the full
algebra~\cite{vilainkautz1986,nebel1995}; the library says so at the
call site rather than letting a clean pass be read as a proof.

This is rarely the right surface for application code, however ---
business questions tend to name unions of Allen relations.
``Does the meeting fall within the workday?'' is the set
$\{\textit{equals}, \textit{starts}, \textit{during}, \textit{finishes}\}$;
``does this period touch that one?'' is the complement of
$\{\textit{precedes}, \textit{preceded\_by}\}$. Hand-rolled inline
matching on relation lists is verbose and obscures the intent, so the
library promotes the common unions to first-class predicates ---
\texttt{within?/2}, \texttt{overlaps?/2}, \texttt{adjacent?/2},
\texttt{contains?/2} --- each defined as a specific subset of Allen's
relations, held in one table, and each readable aloud as the English
question it answers.

The predicate vocabulary is not decorative. It is a design discipline:
when an example in the cookbook reaches for a relation list or a
manual endpoint comparison, that is taken as evidence that the
abstraction is missing, and the predicate is added before the example
is written. The same discipline applies to durations
(\texttt{at\_least?}, \texttt{shorter\_than?}) and to interval-set
queries (\texttt{covered?}, \texttt{disjoint?}). The library's
public surface is, deliberately, an English-prose description of
Allen's algebra over ISO~8601 values.

\section{Position among temporal formalisms}

The choice to make intervals primitive places Tempo on one side of a
long-standing divide. Vilain and Kautz~\cite{vilainkautz1986} take
\emph{points} as primitive: an interval is an ordered pair of endpoint
points, and an interval relation reduces to a conjunction of
point-relations $\{<,=,>\}$ among the four endpoints. The reward is
tractability --- reasoning in the point algebra is polynomial, whereas
constraint reasoning over the full interval algebra is intractable in
general~\cite{vilainkautz1986,nebel1995}. Allen and
Hayes~\cite{allenhayes1989} take the opposite stance, axiomatising the
algebra from the single primitive \textit{meets} and reconstructing
points only as derived entities --- \emph{moments} (intervals with no
proper subinterval) and the meeting-places between intervals.
Gr\"uninger and Li's $T_{bounded\_meeting}$~\cite{gruninger2017allenontology}
continues this interval-only lineage, excluding degenerate intervals
and points from the domain entirely; this is the ontology Tempo
adopts.

Tempo's contribution is best read as a synthesis across this divide.
\emph{Ontologically} it is interval-only: there are no instants,
degenerate intervals are excluded, and the half-open seam realises
Hayes' open-connected model rather than Allen's closed intervals.
\emph{Computationally}, however, every end is an integer position on
one frame, and the engine orders spans by their ends --- precisely the
Vilain--Kautz reduction, used as an engine beneath an interval
ontology. Two consequences follow. First, because Tempo's spans are
always \emph{fully grounded} (concrete ends, never disjunctive
qualitative constraints), it never enters the NP-hard regime that
motivated the point algebra: relation queries are four integer
comparisons rather than constraint propagation, and every set
operation is one sweep over sorted ends. Calendars and zones are
consulted only to compute positions, once per end. Second, the
``instant'' is recovered not as a primitive point but as the interval
of one unit at the value's \emph{finest declared resolution} --- a
\emph{moment} in Allen and Hayes' sense, but relativised to
representational precision: a second-resolution datetime is a
one-second interval, a microsecond value a one-microsecond interval.
This resolution-indexed atomicity, tying interval granularity to
ISO~8601 syntax rather than to an absolute atom, and the reading of
written ends by the class of their unit, are the two elements of the
model without a direct precedent in the formalisms it draws on.

\section{Set operations on spans and sets of them}

Closing the library under set-theoretic operations is where the model
earns its keep. The intuition is straightforward --- any two
expressions combine into a (possibly disjoint, possibly empty) set of
spans --- but the engineering requires one answer to the questions of
resolution, of frames, of calendars and of what a result \emph{is}
when its operands were events with identities of their own.

\paragraph*{One sweep.} Every instant-level operation is the same
sweep over the operands' sorted ends with a different question of how
many members of each operand cover a point: \texttt{intersection/3}
keeps a point covered by both, \texttt{difference/3} a point covered by
the first and not the second, \texttt{symmetric\_difference/3} a point
covered by exactly one, and the $n$-ary ``at least $k$ members'' query
a point covered by $k$ or more. A result begins where the answer turns
true and ends where it turns false, so a result's ends are always
operands' ends, and are written back as those operands wrote them, at
their resolution: the difference of a month and a day ends on a day.
Every operation returns a set; a set that happens to have one member
is still a set, which keeps the operations closed and composable
without a case split at the call site.

\paragraph*{Members and covered time.} The textbook operations answer
about \emph{covered time}. Application code as often asks about
\emph{members}: which meetings hit this window, which workdays are not
holidays. Tempo keeps the two apart. \texttt{union/3} preserves
members (two meetings that cover the same hour stay two meetings, each
with its metadata) and the coalesced form is one explicit operation
away; the other operations trim to covered time, and each result is
credited to the operand members that cover it, so that a holiday's
name survives the difference that removed the weekend from it. Each
instant-level operation has a member-preserving companion
(\texttt{members\_overlapping/3}, \texttt{members\_outside/3},
\texttt{members\_in\_exactly\_one/3}), and the canonical form of a
set --- sorted by start, coalesced where members touch --- is one
operation, \texttt{coalesce/2}, which realises the Sum Axiom of
$T_{bounded\_meeting}$ (axiom~15 in~\cite{gruninger2017allenontology}):
three chain-meeting intervals
$\textit{meets}(i,j) \land \textit{meets}(j,k) \land \textit{meets}(k,m)$
are summed into the single interval that runs from $i$'s start to
$m$'s end.

\paragraph*{One expression, one calendar, one zone.} The calendar and
the zone are properties of a whole expression, written once after it
(\texttt{2026-06-15T10:30/T12:00[Europe/Paris][u-ca=hebrew]}), never of
one end or one member: an interval whose two ends are written in two
calendars, or two named zones, is refused at construction. Across
expressions the rule is explicit conversion. A set operation brings its
second operand to the calendar of its first; two zoned operands meet on
the UTC line and the result is written in the first operand's zone; a
floating operand beside a zoned one is a named error unless the caller
places it. The result's calendar and zone are therefore the first
operand's, which preserves the caller's intent for display, and what
cannot be written in one calendar and one zone is an error rather than
another value. An earlier version of this design projected every
operand to a common line and let the first operand win silently; the
stricter rule adopted the alternative --- refuse heterogeneous frames,
convert calendars explicitly --- that the authors had invited the
symposium to weigh, and it is one of the points the talk would revisit.

\paragraph*{Cross-calendar positions.} A date is placed on the day
scale as a proleptic Gregorian day number, the fixed date of Reingold
and Dershowitz, which every calendar converts to and from without a
path through another calendar; a reading of the clock is placed as a
microsecond tick on the UTC line through the zone. The word ``instant''
is avoided because Gr\"uninger and Li's ontology excludes points from
its domain; the projection is a computational device, an integer used
to order ends on one frame.

\paragraph*{Daylight-saving transitions.} Because every span is read
through its zone to positions on the UTC line, transitions fall out of
the walk rather than being handled by hand. A day in Paris on which the
clock springs forward has twenty-three hours because its two ends are
twenty-three hours apart; an hour the clock skips in part is as long as
the clock shows it; a day a zone leaves out (Samoa's 30 December 2011)
is no day of a walk through that December. Two cases are settled at
construction, so that set operations on positions need no further
disambiguation. A wall time that falls in a DST \emph{gap} (the
non-existent hour at spring-forward) is refused as a value: a
\texttt{Tempo.ZoneError} whose reason is \texttt{:gap}, which is a
stronger guarantee than addressing the bug, because the malformed value
cannot be constructed. A wall time that falls in a DST \emph{fold} (the
doubled hour at fall-back) is disambiguated by an explicit offset in
the IXDTF suffix, per RFC~9557~\S 4.5: the value
\texttt{2024-11-03T01:30:00-04:00[America/New\_York]} unambiguously
picks the pre-fall-back (EDT) showing and
\texttt{2024-11-03T01:30:00-05:00[America/New\_York]} picks the
post-fall-back (EST) showing. The disambiguator is part of the string
representation and round-trips through serialisation --- a property
that distinguishes this design from meta-attribute schemes (such as
Python's \texttt{fold} parameter) in which the choice is carried
out-of-band and lost on serialisation.

\paragraph*{Expressions with no end.} A recurrence with no end is a set
with infinitely many members, and the operations answer it exactly or
not at all: an intersection with a bounded operand lies within that
operand and is taken there; a union of two endless operands is an
endless set in time order, read lazily; an intersection of two endless
operands, which may be searched for ever, is refused unless the caller
gives a window.

\section{Implementation}
\label{sec:implementation}

Tempo is a working library, and the model above is its second
implementation. The first grew to 82,875 lines by accretion and carried
the same handful of concepts in parallel forms: a review in October
2026 counted ten walks, six sweeps, five copies of the Allen table,
four selection engines and two arithmetics. The second was written to
the nine laws in one pass and is 26,862 lines, of which 14,524 are
code and 7,806 are the engine (positions and frames, the calendar and
zone contracts, expansion, the walk, the sweep, the classifier,
selection, arithmetic). Each concept is implemented once, and the
calendar is reached through one contract of eighteen questions that
never asks what kind of calendar it is talking to, so that the
Gregorian, Hebrew, Islamic, Chinese, Julian and fiscal calendars join
by the same door.

The implementation is held to oracles written apart from it: an
independent expander of 162 recurrence rules from three starts, itself
held to RFC~5545's own examples; a check of every set operation on
every pair of sets of up to two members between five points by integer
arithmetic on the standard library's dates; Allen's composition table
held to what three spans give; twelve zones over twenty-one days each
held to the standard library's clock; and a totality check that every
operation on every expression of a corpus is an answer or a named
error. Against the first implementation, reading text is thirty times
faster, relating two values two to three times slower, and expanding a
recurrence rule two to six times slower, the price of one selection
engine where the first had a fast path for weekdays; a positioned
representation through the engine and a compiled matcher for plain
rules in front of it are the next steps and change nothing of the
model.

\section{Discussion and open questions}
\label{sec:discussion}

This submission's contribution is a
\emph{design} contribution rather than a \emph{result} contribution,
and the talk would aim to make the design choices legible to an
audience whose research is the theory the library implements. Four
questions seem particularly worth the TIME audience's attention:

\begin{enumerate}
  \item \textbf{Written ends.} The rule of Section~\ref{sec:model},
    that an end written to a calendar unit is included and an end
    written to a clock unit is a boundary, reconciles the inclusive
    reading of EDTF and ISO~8601-2 Annex~A with a half-open model by
    appeal to how the two kinds of end are said aloud. It is exact
    and it round-trips, but it makes \texttt{2026-06-15T10/2026-06-16}
    thirty-eight hours where a reader used to boundaries expects
    fourteen. Is a rule keyed to the class of the unit the right
    reconciliation, or should the standard carry a notation for an
    exclusive end?

  \item \textbf{Granularity and resolution.} Tempo's
    one-unit-of-the-written-resolution rule for values is
    operationally simple but does not reference the more general
    treatment of temporal granularity in the database
    literature~\cite{bettini-granularities} nor the calendric type
    constraints in CaTTS~\cite{bry-spranger-multicalendar}. Is the
    simple rule a sufficient model in practice, or does it accumulate
    edge cases as the calendar set grows (lunar, sidereal, fiscal)?
    The three frames, and the cycle in particular, are the places the
    question is felt.

  \item \textbf{Allen's algebra over sets.} The thirteen
    relations, and the $T_{bounded\_meeting}$ ontology that
    underlies them, are defined on convex intervals
    only~\cite{gruninger2017allenontology}. Their lifting to sets
    admits several reasonable definitions, and the library now offers
    three: the covered-time lifting through the predicates
    (\emph{within} is an empty difference, \emph{overlaps} a
    non-empty intersection), the existential lifting through the
    member-preserving companions, and the set-of-relations lifting
    through \texttt{relation\_matrix/2}, which returns the relation
    holding between every pair of members. Which should be the
    \emph{default} is not obvious, and a principled extension of the
    ontology to non-convex sets would settle it.

  \item \textbf{Cross-calendar comparison.} Comparing a Hebrew date
    with a Gregorian one resolves cleanly when both are placed on the
    same day scale, but the scale is itself a modelling choice (TAI?
    UT1? a calendar-specific civil time?). Tempo's choice --- the
    proleptic Gregorian day number for dates and UTC for readings of
    the clock, with calendars converted explicitly and frames never
    mixed --- is pragmatic; a principled framework would be a
    worthwhile contribution from the symposium back to the library.
\end{enumerate}

The wider goal is to demonstrate, by existence proof, that the
theoretical structure of temporal reasoning translates into an
ordinary application library --- one that an engineer writes
\texttt{Tempo.overlaps?(a, b)} into without thinking of it as a
research artefact. The same primitive reaches beyond set algebra: a
thin constraint layer reduces a network of partially known intervals to
a Simple Temporal Problem~\cite{dechter1991}, reproducing the
chronological-network model of Levy et al.~\cite{levy2020} --- realised
in the ChronoLog tool~\cite{chronolog} --- as a library feature rather
than a bespoke application.\footnote{Developed in a companion paper,
  \emph{Chronological Networks as a Library Primitive: ChronoLog over an
  Interval-Native Date/Time Stack}.} The talk would walk through the
design, the choices made, the choices still open, and invite the
symposium to sharpen them.

\begin{thebibliography}{18}

\bibitem{allen1983}
  James F. Allen.
  \newblock Maintaining knowledge about temporal intervals.
  \newblock \emph{Communications of the ACM}, 26(11):832--843, 1983.
  \newblock \doi{10.1145/182.358434}

\bibitem{vilainkautz1986}
  Marc B. Vilain and Henry A. Kautz.
  \newblock Constraint propagation algorithms for temporal reasoning.
  \newblock In \emph{Proceedings of the 5th National Conference on
  Artificial Intelligence (AAAI-86)}, pages 377--382, 1986.

\bibitem{allenhayes1989}
  James F. Allen and Patrick J. Hayes.
  \newblock Moments and points in an interval-based temporal logic.
  \newblock \emph{Computational Intelligence}, 5(3):225--238, 1989.
  \newblock \doi{10.1111/j.1467-8640.1989.tb00329.x}

\bibitem{nebel1995}
  Bernhard Nebel and Hans-J\"urgen B\"urckert.
  \newblock Reasoning about temporal relations: A maximal tractable
  subclass of Allen's interval algebra.
  \newblock \emph{Journal of the ACM}, 42(1):43--66, 1995.
  \newblock \doi{10.1145/200836.200848}

\bibitem{iso8601-1}
  International Organization for Standardization.
  \newblock \emph{ISO 8601-1:2019 --- Date and time --- Representations
  for information interchange --- Part 1: Basic rules}, 2019.

\bibitem{iso8601-2}
  International Organization for Standardization.
  \newblock \emph{ISO 8601-2:2019 --- Date and time --- Representations
  for information interchange --- Part 2: Extensions}, 2019.

\bibitem{ixdtf}
  Ujjwal Sharma and Carsten Bormann.
  \newblock \emph{Date and Time on the Internet: Timestamps with
  Additional Information}.
  \newblock RFC 9557, IETF, April 2024.
  \newblock \doi{10.17487/RFC9557}

\bibitem{edtf}
  Library of Congress.
  \newblock \emph{Extended Date/Time Format (EDTF) Specification}, 2019.
  \newblock \url{https://www.loc.gov/standards/datetime/}

\bibitem{bettini-granularities}
  Claudio Bettini, Sushil Jajodia, and Sean X. Wang.
  \newblock \emph{Time Granularities in Databases, Data Mining, and
  Temporal Reasoning}.
  \newblock Springer, 2000.

\bibitem{tiwari2025time}
  Shrey Tiwari, Serena Chen, Alexander Joukov, Peter Vandervelde,
  Ao Li, and Rohan Padhye.
  \newblock It's about time: An empirical study of date and time bugs
  in open-source Python software.
  \newblock In \emph{Proceedings of the 22nd International Conference
  on Mining Software Repositories (MSR '25)}. IEEE/ACM, 2025.
  \newblock Distinguished Paper Award.
  \newblock \url{https://rohan.padhye.org/files/datetimebugs-msr25.pdf}

\bibitem{liva2018}
  Giovanni Liva.
  \newblock \emph{Modeling Time in Java Programs for Automatic Error
  Detection}.
  \newblock PhD thesis, Alpen-Adria-Universit\"at Klagenfurt, 2018.

\bibitem{bry-spranger-multicalendar}
  Fran\c{c}ois Bry and Stephanie Spranger.
  \newblock Towards a multi-calendar temporal type system for
  (Semantic) Web query languages.
  \newblock In \emph{Principles and Practice of Semantic Web
  Reasoning (PPSWR 2004)}, volume 3208 of \emph{Lecture Notes in
  Computer Science}, pages 90--104. Springer, 2004.
  \newblock \doi{10.1007/978-3-540-30122-6\_8}

\bibitem{gruninger2017allenontology}
  Michael Gr\"uninger and Zhuojun Li.
  \newblock The time ontology of Allen's interval algebra.
  \newblock In \emph{24th International Symposium on Temporal
  Representation and Reasoning (TIME 2017)}, volume 90 of
  \emph{LIPIcs}, pages 16:1--16:16. Schloss
  Dagstuhl -- Leibniz-Zentrum f\"ur Informatik, 2017.
  \newblock \doi{10.4230/LIPIcs.TIME.2017.16}

\bibitem{dyreson2021logsegmented}
  Curtis E. Dyreson and M. A. Manazir Ahsan.
  \newblock Achieving a sequenced, relational query language with
  log-segmented timestamps.
  \newblock In \emph{28th International Symposium on Temporal
  Representation and Reasoning (TIME 2021)}, volume 206 of
  \emph{LIPIcs}, pages 14:1--14:13. Schloss
  Dagstuhl -- Leibniz-Zentrum f\"ur Informatik, 2021.
  \newblock \doi{10.4230/LIPIcs.TIME.2021.14}

\bibitem{dyreson2023logsegmented}
  Curtis E. Dyreson.
  \newblock Optimization of nonsequenced queries using log-segmented
  timestamps.
  \newblock In \emph{30th International Symposium on Temporal
  Representation and Reasoning (TIME 2023)}, volume 278 of
  \emph{LIPIcs}, pages 13:1--13:15. Schloss
  Dagstuhl -- Leibniz-Zentrum f\"ur Informatik, 2023.
  \newblock \doi{10.4230/LIPIcs.TIME.2023.13}

\bibitem{dechter1991}
  Rina Dechter, Itay Meiri, and Judea Pearl.
  \newblock Temporal constraint networks.
  \newblock \emph{Artificial Intelligence}, 49(1--3):61--95, 1991.
  \newblock \doi{10.1016/0004-3702(91)90006-6}

\bibitem{levy2020}
  Eythan Levy, Gilles Geeraerts, Fr\'ed\'eric Pluquet, Eli Piasetzky,
  and Alexander Fantalkin.
  \newblock Chronological networks in archaeology: A formalised scheme.
  \newblock \emph{Journal of Archaeological Science}, 124:105225, 2020.
  \newblock \doi{10.1016/j.jas.2020.105225}

\bibitem{chronolog}
  ChronoLog.
  \newblock \url{http://chrono.ulb.be}.

\end{thebibliography}


\end{document}
